\begin{abstract}
Large language model (LLM) benchmark scores are routinely compared across model generations, yet whether individual items measure the same construct over time is rarely tested. We apply differential item functioning (DIF) analysis from educational measurement to the METABENCH response matrix (5,227 models on MMLU, replicated across six benchmarks), validated through semi-synthetic injection with a cross-subject matching guarantee. Between 2023 and 2024 model cohorts, 18--31\% of items exhibit large temporal DIF, 45--95$\times$ above random baselines, while aggregate rankings remain near-perfect ($\rho \geq 0.994$). DIF flags are near-independent of web-overlap contamination labels ($\phi = 0.012$), indicating that global exposure and differential functioning are complementary diagnostic dimensions. Base and instruct models exhibit opposite DIF directions, confirming structured capability shifts beyond contamination. We propose a three-step Monitor/Triage/Correct audit protocol with empirically grounded thresholds for benchmark maintainers. LLM benchmarks need psychometric monitoring, not just score reporting\footnote{Code and data are available at \url{https://anonymous.4open.science/r/irt_equated_bench_regen-2431/}.}.
\end{abstract}
